\documentclass{jsarticle}
\usepackage[dvipdfmx]{graphicx}
\title{レポートタイトル}

\author{学生番号XXX-XXXX アカリク太郎}
\date{\today}
\begin{document}
\maketitle
\section{Cloud LaTeXへようこそ}

Cloud LaTeXは，\LaTeX を使った文書の作成・管理をクラウド上で行えるWebサービスです．
\LaTeX を使うと，複雑な数式
\begin{equation}
  \frac{\pi}{2} =
  \left( \int_{0}^{\infty} \frac{\sin x}{\sqrt{x}} dx \right)^2 =
  \sum_{k=0}^{\infty} \frac{(2k)!}{2^{2k}(k!)^2} \frac{1}{2k+1} =
  \prod_{k=1}^{\infty} \frac{4k^2}{4k^2 - 1}
\end{equation}
を含んだ読みやすくきれいな文書作成ができます．

本サービスは，\LaTeX 文書をリアルタイムに保存・コンパイルし，ユーザーアカウント別に管理します．
そのため，本サービスにログインするだけで，どこからでも作業を再開でき，ファイルを持ち歩く必要はありません．
また，様々な \LaTeX テンプレートが用意されているので，手軽に文書を作り始めることができます．

\begin{figure}
 \centering
   \includegraphics[width=60mm]{figures/Sample.png}
 \caption{ここにキャプションを挿入します}
 \label{fig:model}
\end{figure}

Cloud LaTeXでは，作成されるPDFそのままのレイアウトで表示するPDFビューモードがあり，コンパイル画面を確認しながら文書を作成することができます（図\ref{fig:model}）
日本語では，pLaTeX / LuaLaTeX / upLaTeX でのコンパイルが可能です．
また，日本語や英語文書作成だけでなく，中国語・ハングルに対応した XeLaTeX のコンパイルも可能です．
ぜひ使ってみてください．
\end{document}
